\section{Results}

Our validation attempt produced one artifact---HumanEval + Pydantic type extensions---while all hypothesis testing was blocked by infrastructure failures. This section presents the results in three parts: dataset extension success, LLM generation failure, and pipeline architectural validation.

\subsection{Dataset Extension: Artifact Contribution}

The dataset extension pipeline successfully generated 100 Pydantic-annotated prompts from HumanEval problems 0-99, each with function signatures, Pydantic BaseModel schemas, @validator decorators for preconditions, and original test suites.

\textbf{Quantitative Results:}

\begin{table}[h]
\centering
\begin{tabular}{lllll}
\toprule
\textbf{Component} & \textbf{Input} & \textbf{Expected Output} & \textbf{Actual Output} & \textbf{Status} \\
\midrule
DatasetExtender & HumanEval (100) & 100 Pydantic prompts & 100 JSON files & \checkmark PASS \\
\bottomrule
\end{tabular}
\caption{Dataset Extension Results}
\end{table}

Manual inspection of 10 randomly-sampled prompts confirmed type annotations match original signatures, validators encode docstring preconditions, and test cases remain unchanged from HumanEval baseline.

\textbf{Artifact Availability:} Dataset available in \texttt{src/data/humaneval\_pydantic/}, compatible with standard HumanEval evaluation harness. This artifact demonstrates that typed benchmarks can be mechanically generated from existing code generation datasets, addressing a gap in formal verification research.

\subsection{LLM Generation: Infrastructure Failure}

All LLM code generation attempts failed with HTTP 401 authentication errors, blocking hypothesis validation.

\textbf{Quantitative Results:}
\begin{itemize}
\item Total generation attempts: 48 (experiment stopped early after consistent failures)
\item Successful generations: 0 (0\% success rate)
\item Error type: HTTP 401 ``API key is invalid'' (100\% of failures)
\item Retry attempts per request: 3 (exponential backoff: 1s, 2s, 4s)
\item Total retries exhausted: 48 × 3 = 144 failed API calls
\end{itemize}

\textbf{Error Message (Representative):}
\begin{verbatim}
Error code: 401 - {'type': 'error', 'error': {'type':
'authentication_error', 'message': 'API key is invalid.'}}
\end{verbatim}

\subsubsection{Root Cause Analysis}

Three competing explanations were evaluated:

\begin{enumerate}
\item \textbf{Invalid API key:} Error message explicitly states ``API key is invalid''; consistent across all 48 attempts. \textbf{Plausibility: HIGH.}

\item \textbf{API rate limiting:} Would produce HTTP 429 errors, not 401. Anthropic API status showed no outages during experiment window. \textbf{Plausibility: LOW.}

\item \textbf{Network/firewall blocking:} Would produce connection timeouts or DNS errors, not authenticated 401 responses. \textbf{Plausibility: LOW.}
\end{enumerate}

\textbf{Most Likely Interpretation:} Invalid API key (explanation 1). Retry with confirmed-valid key would distinguish between infrastructure failure and hypothesis refutation.

\subsubsection{Implications}

The h-e1 gate criterion---extraction\_success\_rate $\geq$90\%---could not be tested. The hypothesis remains \textbf{theoretically plausible but empirically unvalidated}. All downstream hypotheses (h-m1, h-m2, h-m3) were blocked by this prerequisite failure.

\subsection{Pipeline Architectural Validation}

Despite generation failure, the constraint extraction pipeline was validated on error files (programs containing only error comments instead of valid Python code).

\textbf{Component-Level Results:}

\begin{table}[h]
\centering
\small
\begin{tabular}{lllll}
\toprule
\textbf{Component} & \textbf{Input} & \textbf{Expected} & \textbf{Actual} & \textbf{Status} \\
\midrule
DatasetExtender & HumanEval (100) & 100 prompts & 100 JSON files & \checkmark PASS \\
LLMGenerator & 100 prompts & 100 programs & 48 error files & \texttimes FAIL (infra) \\
PyreExtractor & 48 error files & 0 constraints & 0 constraints & \checkmark PASS (neg) \\
Z3Validator & 0 constraints & 0 quality & 0 quality & \checkmark PASS (neg) \\
MetricsEngine & 0 constraints & rate=0\%, FAIL & rate=0.0\%, FAIL & \checkmark PASS \\
\bottomrule
\end{tabular}
\caption{Pipeline Component Validation}
\end{table}

\textbf{Architectural Soundness:} Pipeline validated on negative cases only (error files processed without crashing; positive case untested). This validates:

\begin{itemize}
\item \textbf{PyreExtractor:} AST parser correctly handles unparseable Python (error comments), returns 0 constraints as expected.
\item \textbf{Z3Validator:} Constraint quality checker correctly identifies 0 quality constraints from empty input.
\item \textbf{MetricsEngine:} Gate evaluation logic correctly computes 0\% extraction rate and triggers FAIL condition.
\end{itemize}

\subsubsection{Implications}

The pipeline is \textbf{architecturally sound and ready for retry} (on negative cases only; positive validation requires actual LLM code). Fixing the API key requires no code changes---only environment configuration. If retry succeeds (extraction\_rate $\geq$90\%), the pipeline will proceed to h-m1 without modification.

\subsection{Visualization Results}

Four figures were generated by the MetricsEngine, though three contain placeholder/zero data due to generation failure:

\textbf{Figure \ref{fig:gate}:} Gate threshold (90.0\%) vs actual extraction rate (0.0\%). Shows clear gate failure.

\textbf{Figure \ref{fig:failure}:} Categorization of extraction failures. 100\% API authentication errors (48/48 attempts). \textbf{This is the only figure with meaningful data.}

\textbf{Figure \ref{fig:complexity}:} Planned analysis of extraction success by program complexity (LOC bins). All bins show 0\% due to no successful generations. Contains placeholder data.

\subsection{Hypothesis Validation Summary}

\begin{table}[h]
\centering
\small
\begin{tabular}{lllllp{4cm}}
\toprule
\textbf{Hyp} & \textbf{Gate} & \textbf{Threshold} & \textbf{Actual} & \textbf{Status} & \textbf{Reason} \\
\midrule
h-e1 & MUST & rate $\geq$90\% & 0.0\% & \texttimes FAIL & Infrastructure block \\
h-m1 & MUST & All extract & N/A & NOT\_STARTED & Prereq h-e1 failed \\
h-m2 & SHOULD & cone <30\%, detect=100\% & N/A & NOT\_STARTED & Prereq failed \\
h-m3 & MUST & median $\geq$2x & N/A & NOT\_STARTED & Prereq failed \\
\bottomrule
\end{tabular}
\caption{Hypothesis Validation Summary}
\end{table}

\textbf{Predictions:}
\begin{itemize}
\item P1 (2-5x speedup): \textbf{UNTESTED} (h-m3 not reached)
\item P2 (speedup scales with size): \textbf{UNTESTED} (h-m3 not reached)
\item P3 (soundness maintained): \textbf{UNTESTED} (h-m2 not reached)
\end{itemize}

\textbf{Assumptions:}
\begin{itemize}
\item A1 (LLM code has extractable constraints): \textbf{UNVERIFIED} (h-e1 blocked)
\item A2 (repair locality): \textbf{UNVERIFIED} (h-m2 not reached)
\item A3 (dependency analysis soundness): \textbf{UNVERIFIED} (h-m2 not reached)
\item A4 (extraction overhead small): \textbf{UNVERIFIED} (h-m1 not reached)
\item A5 (no dynamic features): \textbf{PARTIALLY VERIFIED} (prompt constraints applied, not tested)
\end{itemize}

\subsection{Summary}

\textbf{Artifact:} HumanEval + Pydantic dataset extension (100 typed prompts).

\textbf{Unverified:} All hypothesis claims (incremental SMT speedup, constraint extraction feasibility, dependency analysis soundness).

\textbf{Failure Mode:} Infrastructure (invalid API key), not scientific refutation.

\textbf{Next Step:} Retry h-e1 with valid Anthropic API key. If extraction\_rate $\geq$90\%, proceed to h-m1-m3 validation.
